\documentclass[11pt]{article}
\usepackage[T1]{fontenc}
\usepackage{lmodern,amsmath,amssymb,amsthm,mathtools,microtype}
\usepackage[left=.94in,right=.94in,top=.78in,bottom=.82in]{geometry}
\usepackage{xcolor}
\usepackage[colorlinks=true,allcolors=blue!45!black]{hyperref}
\hypersetup{pdftitle={Liouville signs of short differences in sets of primes},
pdfsubject={Research draft on weighted prime pairs and local tuples}}
\newtheorem{theorem}{Theorem}[section]
\newtheorem{corollary}[theorem]{Corollary}
\newtheorem{proposition}[theorem]{Proposition}
\newtheorem{lemma}[theorem]{Lemma}
\newcommand{\E}{\mathbb E}
\newcommand{\Z}{\mathbb Z}
\newcommand{\one}{\mathbf 1}
\newcommand{\Sf}{\mathfrak S}
\newcommand{\PP}{\mathbb P}
\newcommand{\range}[1]{[\,#1\,]}
\newcommand{\supp}{\operatorname{supp}}
\newcommand{\op}{\mathrm{op}}
\newcommand{\HS}{\mathrm{HS}}
\newcommand{\e}{\mathrm e}
\newcommand{\cD}{\mathcal D}
\newcommand{\cA}{\mathcal A}
\newcommand{\NN}{\mathcal N}
\allowdisplaybreaks[1]
\raggedbottom
\setlength\emergencystretch{1.8em}
\title{\vspace{-1.1em}\bfseries Liouville signs of short differences\\in sets of primes}
\author{\small \href{mailto:admin@cosmosgeek.org}{admin@cosmosgeek.org}\qquad \href{mailto:b321014@iiit-bh.ac.in}{b321014@iiit-bh.ac.in}}
\date{\small Working paper \enspace$\cdot$\enspace November 2026}
\begin{document}
\maketitle
\begin{abstract}
Let \(A\) be a subset of the primes between \(X\) and \(2X\) of fixed positive relative density.
Among pairs \(p<q\) of elements of \(A\) with \(q-p\le 2R\), the integer
\((q-p)/2\) has an even number of prime factors with limiting frequency \(1/2\),
counting multiplicity, whenever \(R/\log X\to\infty\) and \(R\le X/100\).
The estimate is uniform in \(A\). At \(R=\lfloor4\log X\rfloor\) a positive-relative-density
selection can give every such pair the same sign. For fixed-size ordered tuples
in a translated window, all assignments of signs to their pairwise half-differences
have equal limiting frequencies under the window count. We prove a weighted
form for one or two polynomially rough translates by combining a smooth sieve
correlation estimate with Fourier cancellation on the difference variable.
For two rough translates we also construct single-sign selections below the squared-logarithmic scale.
\end{abstract}

\section{Short prime differences}
Write \(\lambda(d)=(-1)^{\Omega(d)}\), where \(\Omega\) counts prime factors
with multiplicity, and set \(\lambda(1)=1\). All logarithms are natural.
For \(A\subset\PP\cap(X,2X]\) and \(\sigma=\pm1\), let \(N_\sigma(A,R)\)
count pairs \(p<q\) in \(A\) satisfying \(q-p\le2R\) and
\(\lambda((q-p)/2)=\sigma\). Write \(N=N_++N_-\), and \(L=\log X\).

\begin{theorem}\label{th:prime}
Suppose \(R\) is a positive integer, \(R/L\to\infty\), and \(R\le X/100\).
For each fixed \(\delta>0\), uniformly over \(A\) with \(|A|\ge\delta X/L\),
\[
    N(A,R)\asymp_\delta XR/L^2,\qquad
    N_\sigma(A,R)=\bigl(\tfrac12+o_\delta(1)\bigr)N(A,R).
\]
For \(R=\lfloor4L\rfloor\) and all sufficiently large \(X\), there is a set
\(A_X\) of primes in the same interval, \(|A_X|\gg X/L\), for which
\(N(A_X,R)\gg X/L\) and one of \(N_+(A_X,R),N_-(A_X,R)\) vanishes.
\end{theorem}

At \(R\) just above the mean prime spacing, a particular interval can have very
few primes. The theorem sums over the dyad and permits the set to depend on
\(X\) and \(R\). Each pair has unit weight in Theorem~\ref{th:prime}, including
pairs with intervening primes. The assertions concern averaged separations.
They give no estimate for the count at the single separation \(q-p=2\).

Rivat, S\'ark\"ozy and Stewart \cite[Theorem 1]{rss} study \(\Omega(a+b)\)
modulo an integer for arbitrary sets in \([1,N]\). Their error has size
\((|A||B|)^{1/2}N/(\log N)^E\). At a local scale \(N\asymp R\) near
\(\log X\), such a saving in \(\log R\) does not compensate for the density
\(1/\log X\) at each prime endpoint. Banks and Shparlinski
\cite[\S7.1]{banks} treat M\"obius-weighted sums of two prime variables,
allowing injective transformations of the summation variables. Their bound
in that application requires the endpoint scale to be at most
\(\exp((\log R)^C)\) for fixed \(C\), when the transformed sum has size at
most \(R\). Lichtman's averaged shifted-prime theorem
\cite[Theorem 1.1]{lichtman} places the M\"obius factor on \(p+h\).
Our oscillating factor is evaluated on the short difference of two selected
vertices.

Tao and Ziegler \cite{tz} and Shao \cite{shao} establish narrow arithmetic
progressions in dense subsets of primes. Shao's Proposition 2.6 gives smooth
majorant correlations at widths exceeding a system-dependent power of \(\log X\)
by an arbitrarily slowly growing factor. His collision analysis and averaged
singular factors are particularly relevant here. We use the Goldston--Y{\i}ld{\i}r{\i}m
estimate in Green and Tao \cite[Appendix D, Theorem D.3]{gt}. The contribution
here is the short-distance sign statistic uniformly over endpoint selections,
its joint version for all edges of a local tuple, and the mean-spacing
counterexample in Theorem~\ref{th:prime}. The tuple indices vary independently.
No arithmetic-progression constraint is imposed on them.

\section{Weighted pairs and tuples}
Fix a positive integer \(s\), a real number \(0<\beta<1\), and one of
\(\mathcal H=\{0\},\{0,2\}\). Put \(j=|\mathcal H|\). Call \(n\) admissible
if \(X<n\le2X\) and every \(n+h\), \(h\in\mathcal H\), has all its prime
factors greater than \(X^\beta\). A vertex weight is any complex function
supported on the admissible integers with absolute value at most \(L^j\).
Throughout this section let \(R\) be a positive integer and assume
\begin{equation}\label{eq:width}
       R/L^j\longrightarrow\infty,\qquad 1\le R\le X/(100s).
\end{equation}
The notation \([R]\) means \(\{1,\ldots,R\}\).

\begin{theorem}\label{th:weighted}
For any two vertex weights \(F,G\), any interval \(I\subseteq[R]\), and
\(a=\lambda\) or the M\"obius function \(\mu\),
\begin{equation}\label{eq:pair}
       \frac{1}{XR}\sum_{d\in I} a(d)\sum_{n\in\Z}F(n)G(n+2sd)=o(1).
\end{equation}
Fix \(k\ge2\) and a nonempty simple graph \(E\) on \(\{1,\ldots,k\}\).
With \(F_1,\ldots,F_k\) any vertex weights and
\(\mathcal X=\Z\cap[X/2,3X]\), one has
\begin{equation}\label{eq:graph}
 \frac{1}{XR^k}\sum_{x\in\mathcal X}
 \sum_{\substack{u_1,\ldots,u_k\in[R]\\\text{pairwise distinct}}}
 \ \prod_{i=1}^k F_i(x+2su_i)
 \prod_{\{i,l\}\in E}\lambda(|u_i-u_l|)=o(1).
\end{equation}
The errors are uniform in the weights and in \(I\), for the indicated fixed
parameters and every sequence satisfying \eqref{eq:width}.
\end{theorem}

\begin{corollary}\label{co:tuple}
Let \(A\) be a set of admissible integers of cardinality at least
\(\delta X/L^j\). Count ordered tuples by summing over \(x\in\mathcal X\)
and distinct \(u_i\in[R]\) for which every \(x+2su_i\) belongs to \(A\).
If \(T\) is their number, then \(T\asymp_{\delta,k,s,\beta} XR^k/L^{jk}\).
For each prescribed array of signs
\(\sigma_{il}\), \(1\le i<l\le k\), the proportion with
\(\lambda(|u_i-u_l|)=\sigma_{il}\) for all these pairs is
\(2^{-\binom{k}{2}}+o_\delta(1)\), with \(k,s,\beta\) fixed.
\end{corollary}

This sampling counts a physical tuple once for every window containing it.
For \(s=1\), an ordered tuple of odd primes of diameter \(2d<2R\)
has \(R-d\) anchors. Thus the tuple theorem specifies its sampling even
when \(R\) is a substantial fraction of \(X\). The two-translate case gives a
relative frequency when the selected rows have the stated density.
Its signed estimate \eqref{eq:pair} is absolute and needs no lower count.

\section{Moments of the sieve weight}
Choose a real smooth function \(\chi\) supported in \([-1,1]\) with
\(\chi(0)=1\). For this proof take \(N_0=\lceil4X\rceil\),
\(D=N_0^\eta<X^\beta\) with a fixed \(\eta>0\) sufficiently small for all
the finitely many applications below. Set
\[
 S(n)=\sum_{d\mid n}\mu(d)\chi(\log d/\log D),\qquad
 \nu_0(n)=(\log D)S(n)^2,\qquad
 \nu(n)=C_0\prod_{h\in\mathcal H}\nu_0(n+h).
\]
Take \(C_0\) fixed and large enough that \(C_0(\log D)^j\ge L^j\).
All arguments of these weights in the proof are positive; on an admissible
row all the divisor sums in its majorant equal one. Write \(\nu_t=\nu(x+2st)\).
We use normalized counting measure for \(\E\).

Here is the precise consequence of \cite[Theorem D.3]{gt} used in the proof.
For distinct integer shifts \(t_1,\ldots,t_m\) with \(\max_i|t_i|=O(X)\),
positive fixed integers \(a_i\), and an \(x\)-interval of length \(\asymp X\)
inside \([-N_0,N_0]\) where the shifted arguments are positive, put
\[
 \Delta=\prod_{i<l}|t_i-t_l|,\qquad
 \Sf(\mathbf t)=\prod_p
 \frac{1-|\{t_1,\ldots,t_m\}\bmod p|/p}{(1-1/p)^m}.
\]
An empty product is one. Here \(c_{\chi,a}\) denotes the sieve factor in
\cite[Definition D.1]{gt}; it depends only on the fixed function \(\chi\)
and the fixed integer \(a\). In particular,
\(c_{\chi,2}=\int_0^\infty |\chi'(t)|^2\,dt\), by
\cite[Lemma D.2]{gt}. With this normalization, the average over the
interval satisfies
\begin{align}
 \E_x \prod_{i=1}^m [(\log D) S(x+t_i)^{a_i}]
 =\Big(\prod_{i=1}^m c_{\chi,a_i}\Big)\Sf(\mathbf t)
 +O\!\left(L^{-1/20}
       \exp\!\left(C\sum_{p\mid\Delta}p^{-1/2}\right)\right).
 \label{eq:GT}
\end{align}
The interval-length normalization in place of its Euclidean volume costs
\(O(X^{-1})\) times the main coefficient and is absorbed here.
The forms \(x+t_i\) have bounded scaled coefficients. Two become proportional
over \(\mathbb F_p\) precisely when their shifts coincide modulo \(p\).
Grouping equal integer shifts before using \eqref{eq:GT} is essential.

\begin{lemma}\label{lem:mom}
Let \(b\) be fixed and the \(b\) indices be averaged independently on
\([R]\) or \([2R]\), with \(x\) uniform on \(\mathcal X\). Form the \(jb\)
shifts \(2st_i+h\), \(i\le b,\ h\in\mathcal H\).
Under \eqref{eq:width},
\(\E\prod_{i=1}^b\nu_{t_i}=O_b(1)\), with the fixed parameters understood.
The contribution where two shifts coincide as integers is
\[
       O_b\!\left(\sum_{c=1}^{b-1}(L^j/R)^{b-c}\right).
\]
On the remaining tuples the \(x\)-average equals
\(C_0^b c_{\chi,2}^{jb}\Sf(\mathbf t')\), with average absolute error
\(O_b(L^{-1/20})\), where \(\mathbf t'\) lists the \(jb\) shifts.
If \(\Sf_{\le w}\) is the product restricted to \(p\le w\), then
\[
 \lim_{w\to\infty}\limsup_{X\to\infty}
 \E_{\mathbf t}\one_{\rm distinct\ shifts}
       |\Sf(\mathbf t')-\Sf_{\le w}(\mathbf t')|=0.
\]
\end{lemma}

\begin{proof}
Join two index slots if some of their shifts coincide. In a connected
component, equality determines each index from one free index by a
bounded fixed offset; inconsistent equalities give an empty set.
For fixed \(b,s,\mathcal H\) there are finitely many patterns. A pattern
with \(c\) components has \(O_b(R^c)\) index tuples. If there are \(m\)
distinct shifts overall, then \(m\ge jc\): each component contains at
least one full translate of \(\mathcal H\), disjoint from the others.
A shift of multiplicity \(r\) contributes
\(\nu_0(x+t)^r=(\log D)^{r-1}[(\log D)S(x+t)^{2r}]\).
Hence \eqref{eq:GT} carries an extra factor \((\log D)^{jb-m}\).

For \(\Delta\) formed with distinct shifts, all fixed moments of
\(B(\Delta)=\exp(C\sum_{p\mid\Delta}p^{-1/2})\) have bounded average
over the free indices, with normalization \(R^{-c}\), on the specified
pattern. Indeed, differences within a component are fixed nonzero
integers. Each other difference is of the form \(2s(v-v')+a\) with
bounded fixed \(a\). Its nonzero absolute values are \(O_s(R)\), each with
multiplicity \(O_s(R)\) in the two free indices. For every \(C'>0\),
\begin{equation}\label{eq:euler}
 \frac1Y\sum_{n\le Y}\exp\!\Big(C'\sum_{p\mid n}p^{-1/2}\Big)
 \le\prod_p\Big(1+\frac{\exp(C'/\sqrt p)-1}{p}\Big)\ll_{C'}1.
\end{equation}
This follows by opening a nonnegative squarefree divisor sum and bounding
\(\lfloor Y/d\rfloor/Y\le1/d\). H\"older across the bounded list of
differences proves the assertion, including after restrictions on the free
indices. Both the full singular series and every partial product are bounded
above by a constant times \(B(\Delta)\): beyond the fixed small primes the
logarithm of a local factor is \(O_b(p^{-2}+\one_{p\mid\Delta}/p)\).
Zero factors are allowed. Thus each collision pattern contributes at most
\[
  O_b\big(R^{c-b}L^{jb-m}\big)
       \ll_b (L^j/R)^{b-c}.
\]
For \(c=b\) all shifts are distinct and \eqref{eq:GT} gives the stated
main and averaged error directly.

For the tail take \(w>2jb\), also exceeding the prime divisors of the
fixed nonzero within-slot differences. The expectation, on the distinct
tuples, of \(\sum_{p>w,\ p\mid\Delta}1/p\) is
\(O_{b,s}(\sum_{p>w}p^{-2})\). For a nonconstant difference this follows
by the same two-index multiplicity bound and by
\(\sum_{1\le n\le Y,\ p\mid n}1\le Y/p\).
The logarithm of the tail product therefore tends to zero in measure in
the indicated order of limits. The uniform second-moment bounds on both
products give the asserted \(L^1\) convergence by Cauchy--Schwarz on the
exceptional event. A zero small-prime factor makes both products zero.
\end{proof}

\section{A signed edge after two Cauchy--Schwarz inequalities}
We prove a bound usable for both statements of Theorem~\ref{th:weighted}.
Choose \(u\) uniform on \([R]\), and \(v\) uniform on either \([R]\) or
\([2R]\). The auxiliary indices \(\mathbf t\) are finitely many independent
uniform elements of \([R]\). Write \(z=(x,\mathbf t)\) with \(x\) uniform
on \(\mathcal X\), and put \(W(z)=\prod_i\nu_{t_i}\).
For functions \(|f_z(u)|\le\nu_u,\ |g_z(v)|\le\nu_v\) consider
\[
       T=\E_z W(z)\E_{u,v} f_z(u)g_z(v) h(v-u).
\]
We use two kernels, extended by zero elsewhere:
\(h(d)=\one_I(d)a(d)\) for the pair assertion, and
\(h(d)=\one_{1\le|d|\le 2R}\lambda(|d|)\) for the tuple assertion.

Let \(v_1,v_2\) be independent copies of \(v\), and similarly \(u_1,u_2\).
Then
\begin{align}
 |T|^4&\le (\E_{z,u}W\nu_u)^2
       (\E_{z,v_1,v_2}W\nu_{v_1}\nu_{v_2})\, B, \label{eq:cs}\\
 B&=\E_{z,v_1,v_2}W\nu_{v_1}\nu_{v_2}
 \left|\E_u \nu_u h(v_1-u)\overline{h(v_2-u)}\right|^2. \nonumber
\end{align}
Indeed, Cauchy--Schwarz in \((z,u)\) with measure \(W\nu_u\) first gives
\(|T|^2\le(\E W\nu_u)\E W\nu_u|\E_v g_z(v)h(v-u)|^2\).
Open the \(v\)-square, integrate in \(u\), and apply Cauchy--Schwarz with
measure \(W\nu_{v_1}\nu_{v_2}\), using the bounds on \(g\).
The other two factors in \eqref{eq:cs} are bounded by Lemma~\ref{lem:mom}.
The remaining expression is nonnegative and contains no selected weights.
Opening the \(u\)-square leaves the four factors
\begin{equation}\label{eq:cycle}
 h(v_1-u_1)\overline{h(v_2-u_1)}
 \overline{h(v_1-u_2)}h(v_2-u_2)
\end{equation}
and one \(\nu\) for each of the four indices and each auxiliary slot.
Put \(b=4+\#\{\text{auxiliary slots}\}\).

By Lemma~\ref{lem:mom}, tuples with colliding physical shifts contribute
\(o(1)\) to \(B\) in absolute value. The generic \(x\)-average can be
replaced by its stated constant times \(\Sf_{\le w}\), with errors tending
to zero first in \(X\) and then in \(w\). For fixed \(w\) the truncated
product, evaluated by its residue formula with denominator exponent \(jb\),
extends to a bounded periodic function of all indices, with period
\(P_w=\prod_{p\le w}p\). Restoring tuples with a collision in this periodic
average costs \(O_w(R^{-1})\). Expanding by residue tuples leaves products
of separate indicator functions at each index. Auxiliary indices give
bounded scalars. It remains to average \eqref{eq:cycle} against four
one-variable residue and interval masks.

Embed the indices in \(G_M=\Z/M\Z\) with \(M=16R+1\), using normalized
counting measure. Put the kernel on its representatives in \([-2R,2R]\)
and zero elsewhere, and let \(K\) be the operator with kernel \(h(v-u)\).
There is no cyclic overlap on the ranges in use, and changing normalization
costs a bounded factor. Fourier diagonalization gives
\[
 \|K\|_{\op}
 \le\frac1M\sup_{\theta\in\mathbb R}
          \left|\sum_{|d|\le2R} h(d)\e^{2\pi i\theta d}\right|=o(1),
 \qquad \|K\|_{\HS}^2=\E_{d\in G_M}|h(d)|^2\le1.
\]
For \(\mu\), the Fourier estimate follows from Davenport
\cite{dav}, also \cite[(1.1)]{gtnil}. Uniformity in \(I\) follows by
subtracting prefixes: prefixes up to \(\sqrt R\) cost \(O(\sqrt R)\),
and the remaining ones satisfy the uniform logarithmic saving. For
\(\lambda\), use \(\lambda(n)=\sum_{m^2\mid n}\mu(n/m^2)\), which follows
on prime powers and then by multiplicativity. Terms with \(m>B_0\)
cost \(O(R/B_0)\) before division by \(M\); for each \(m\le B_0\) apply
the same estimate with phase \(m^2\theta\). Take \(R\to\infty\) and then
\(B_0\to\infty\). Reflection handles the negative half.

The masked cycle is the trace of a product of four copies of \(K,K^*,K,K^*\)
in cyclic order with multiplication operators for the intervening masks.
Each mask is a contraction. The absolute trace is at most
\(\|K\|_{\op}^{\,2}\|K\|_{\HS}^{\,2}\) by the Hilbert--Schmidt
Cauchy--Schwarz inequality, bounding two factors in operator norm.
For fixed \(w\) there are only finitely many residue tuples.

To make the width condition and the order of limits explicit, write
\(\mathbf r=(u_1,u_2,v_1,v_2,\mathbf t)\). The indices in this vector
have exactly the independent distributions used in the opened square.
Let \(\mathbf r'\) list their \(jb\) physical shifts \(2sr_i+h\),
\(h\in\mathcal H\), and set
\[
 \mathcal E_b(X,R;w)=
 \E_{\mathbf r}\one_{\rm distinct\ physical\ shifts}
       |\Sf(\mathbf r')-\Sf_{\le w}(\mathbf r')|,
\]
where the summand is defined to be zero on collisions. Combining the
preceding estimates with \eqref{eq:cs} gives, for sufficiently large
\(X\) along a sequence satisfying \eqref{eq:width},
\begin{align}
 |T|^4 \le{}& C_b\left\{
       \sum_{c=1}^{b-1}(L^j/R)^{b-c}
       +L^{-1/20}+\mathcal E_b(X,R;w)\right\} \nonumber\\
 &{}+C_{b,w}\left\{R^{-1}
       +\|K\|_{\op}^{\,2}\|K\|_{\HS}^{\,2}\right\}.
 \label{eq:combined-error}
\end{align}
The constants may depend on the fixed parameters
\(b,s,\mathcal H,\beta,\chi,\eta,C_0\), but not on \(X,R,I\) or any
selected weight; \(C_b\) is also independent of \(w\).
Lemma~\ref{lem:mom} states that
\[
 \lim_{w\to\infty}\limsup_{X\to\infty}
       \mathcal E_b(X,R;w)=0
\]
along every such sequence. First hold \(w\) fixed and let \(X\to\infty\):
the collision and sieve errors vanish, while \(R\to\infty\) makes the
Fourier term vanish uniformly in \(I\). Then let \(w\to\infty\).
No rate of decay in \(w\) is asserted or needed. This proves
\(T=o(1)\) uniformly in all the separated weights, even when
\(R/L^j\) tends to infinity arbitrarily slowly.

\begin{proof}[Completion of Theorem~\ref{th:weighted}]
For \eqref{eq:pair}, take no auxiliary slots,
\(f_x(u)=F(x+2su)\), \(g_x(v)=G(x+2sv)\), \(v\in[2R]\).
Every counted \((n,d)\) occurs exactly \(R\) times: choose \(u\in[R]\),
put \(v=u+d\) and \(x=n-2su\). All such \(x\) lie in \(\mathcal X\).
The normalized pair sum is \(2|\mathcal X|T/X\).

For \eqref{eq:graph} choose an edge of \(E\) with endpoints \(u,v\)
and fix the other \(k-2\) indices in \(z\).
All other edges incident to \(u\) or \(v\) give bounded factors of
\(u\) alone or \(v\) alone once \(z\) is fixed. The auxiliary vertex
weights divided by \(W\), and the remaining edges, form a bounded
common factor, absorbed in \(f_z\); set ratios to zero where \(W=0\).
Take \(v\in[R]\). Extend the edge functions harmlessly at zero and
remove repeated indices at the end using the collision estimate of
Lemma~\ref{lem:mom}. This proves the assertion. The same proof allows
arbitrary bounded factors on all untested edges provided each depends
only on \(x\) and its own two endpoint indices.
\end{proof}

\section{Counting tuples and the endpoint example}
\begin{proof}[Proof of Corollary~\ref{co:tuple} and the supercritical part of
Theorem~\ref{th:prime}]
Write \(f=L^j\one_A\). Every \(n\in A\) and \(u\in[R]\) has an anchor in
\(\mathcal X\), so \(\E_{x,u} f(x+2su)=L^j|A|/|\mathcal X|\gg\delta\).
Jensen gives a lower bound \(\gg_k\delta^k\) for the average of
\((\E_u f(x+2su))^k\). Repetitions contribute \(o(1)\), and the unrestricted
upper is \(O(1)\), by Lemma~\ref{lem:mom}. For distinct indices expand
\(2^{-\binom k2}\prod_{i<l}(1+\sigma_{il}\lambda(|u_i-u_l|))\).
Each nonconstant term tends to zero by \eqref{eq:graph}. Dividing by the
unsigned mass proves the corollary.

For primes use \(j=s=1\) and, for example, \(\beta=1/3\).
The two-vertex lower just proved counts each positive-gap pair in both
orders with at most \(R\) anchors per order. It gives
\(N\gg_\delta XR/L^2\).
The exact pair embedding above and the two-slot moment bound give
\(N\ll XR/L^2\), now taking the positive interval kernel bounded by one.
Equation~\eqref{eq:pair} with \(F=G=L\one_A\) gives
\(N_+-N_-=o(XR/L^2)\), proving the two relative frequencies.
\end{proof}

\begin{proof}[Proof of the critical-width part of Theorem~\ref{th:prime}]
Let \(R=\lfloor4L\rfloor\). Partition the dyad's integers into consecutive
complete bins of length \(R\), dropping the last incomplete bin. Write
\(B\sim X/R\) for the number of bins and \(K_i\) for their prime counts.
The prime number theorem gives \(B^{-1}\sum_iK_i=4+o(1)\).
We claim \(B^{-1}\sum_iK_i^2=O(1)\). For two distinct primes within a bin
their separation \(h\) satisfies \(1\le h<R\). The product of their prime
indicators is \(O(\nu_0(n)\nu_0(n+h)/L^2)\).
For each \(h>0\) apply \eqref{eq:GT} to the two distinct shifts \(0,h\).
The average over \(h\le R\) of both the singular factor and the error's
exponential factor is bounded by \eqref{eq:euler}. Consequently
\[
 \sum_i K_i(K_i-1)\ll XR/L^2,\qquad
 B^{-1}\sum_i K_i^2\ll 1+R^2/L^2\ll1 .
\]
This use of distinct shifts requires no supercritical width assumption.
The contribution of bins with \(K_i<2\) to the first moment is at most
one. Cauchy--Schwarz now gives \(\gg B\) bins containing at least two primes.

Choose one pair in each such bin and classify it by its half-difference
sign and the bin index modulo five. One class contains \(M\gg B\) bins.
Let \(A_X\) consist of the two chosen primes in each retained bin.
Two different retained bins have all their integers more than \(2R\)
apart. Each retained intrabin pair has a positive even gap less than \(R\).
Hence the entire count consists of exactly \(M\) pairs of the same sign
and \(|A_X|=2M\gg X/L\). This proves the claim for the same set at both
endpoints. The construction selects primes separately at each \(X\).
\end{proof}


\section{A subcritical obstruction for two rough translates}
The scale in \eqref{eq:width} also has a selection obstruction on the
two-translate rows. The following construction uses actual squarefree
composite integers at both translates.

\begin{proposition}\label{prop:roughsharp}
There exists $\beta_0>0$ with the following property. Fix
$0<\beta<\beta_0$. Let $\mathcal V$ consist of $X<n\le2X$ such that both
$n,n+2$ are squarefree composites all of whose prime factors exceed
$X^\beta$. Join $n,n+2d\in\mathcal V$ for $1\le d\le R$, and give
the edge sign $\lambda(d)$. If $R\to\infty$ and $R=o(L^2)$, then for each
$\sigma\in\{-1,1\}$ one can delete $O_\beta(XR/L^4)$ vertices so that
all remaining edges have sign $\sigma$ and at least
$c_\beta XR/L^4$ edges remain. Here $|\mathcal V|\gg_\beta X/L^2$.
\end{proposition}

\begin{proof}
We give the sieve estimates with their dependence on $\beta$.
For at most six distinct shifts $t_i$ with $|t_i|\ll R$, let
$\rho(p)$ be the number of forbidden residue classes
$-t_i\bmod p$, using just the shifts being sifted at $p$.
At a squarefree modulus $d$, the count on any interval of length
$\asymp X$ is its length times $\rho(d)/d$, with error $O(\rho(d))$,
where $\rho$ is extended multiplicatively. Up to level $D=X^{1/10}$
the sum of these errors is $O(D(1+\log D)^6)$, since
$\rho(d)\le 6^{\Omega(d)}\le \tau_6(d)$ for squarefree $d$.
An inadmissible small prime leaves no survivors. In all other cases
the dimension bound for the sieve is uniform: for $p>12$ use
$\rho(p)\le6$ and Mertens' theorem, and bound the finitely many
smaller factors separately. Choose $b>0$ small enough that the
fundamental lemma at level $D$ gives fixed-factor upper and lower
bounds with largest sifting prime at most $X^b$
\cite[Theorem~6.9]{fi}. In particular these factors are uniform
in $0<\beta<b$ and the shifts.

For $d\le R$ divisible by $30$, the four shifts $0,2,2d,2d+2$
are distinct and admissible. Sifting all of them up to $X^\beta$
gives $\gg X/(\beta L)^4$ starts. For example, take the starts in
$(X,2X-2R]$ to accommodate both endpoints. The lower bound for
the product of local densities follows by taking
$\rho(p)\le4$ for $p>5$ and using admissibility at the smaller
primes. A positive proportion of the multiples $d=30k$ have each
sign, by $\lambda(30k)=\lambda(30)\lambda(k)$ and the estimate
for $\lambda$ at frequency zero above.

To bound the count in which a specified site is prime, sift that
site further, up to $X^b$, while retaining the original threshold
at the other three. This gains the factor
$\prod_{X^\beta<p\le X^b}(1-1/p)\ll\beta/b$.
The four-site product up to $X^\beta$ is at most
\[
 C(\beta L)^{-4}\prod_{p\mid\Delta}(1+C/p),
 \qquad \Delta=\prod_{i<j}|t_i-t_j|.
\]
The product on the right has bounded average as $d$ ranges over
the multiples of $30$: each varying nonzero difference is linear
in $d$ with fixed coefficients, so the divisor expansion and
H\"older estimate in Lemma~\ref{lem:mom} apply.
Consequently, summed over $d$, removing all four possibilities for
a prime site costs $O(XR/(b\beta^3L^4))$. Removing squares of
primes $p>X^\beta$ costs $O(R(X^{1-\beta}+\sqrt X))$ by summing
$X/p^2+1$ for $p\le\sqrt{3X}$. Choose $\beta_0$ sufficiently small
compared with $b$. Both signs then occur on $\gg_\beta XR/L^4$
edges of $\mathcal V$. The two-site version of the same lower
sieve and the same prime and square removals give
$|\mathcal V|\gg_\beta X/L^2$.

Every start in $\mathcal V$ is $2\bmod3$, so every feasible $d$
is divisible by $3$. The four-site upper sieve, averaged over these
$d$, bounds the total edge count by $O_\beta(XR/L^4)$. Two distinct
edges at a common vertex have signed displacements $2d,2e$, with
$d,e$ distinct nonzero multiples of $3$ and $|d|,|e|\le R$.
The six corresponding shifts $0,2,2d,2d+2,2e,2e+2$ are distinct.
The six-site upper sieve gives $O_\beta(XR^2/L^6)$ incident edge
pairs after averaging over $d,e$. Indeed the nonconstant differences
depend linearly on $d,e$, or $d-e$; each nonzero value occurs
$O(R)$ times among the $O(R^2)$ choices. The same divisor expansion
bounds the averaged singular factors.

Delete both endpoints of every edge of sign $-\sigma$. Each lost
edge of sign $\sigma$ meets such an edge. The number lost is
$O_\beta(XR^2/L^6)=o_\beta(XR/L^4)$, which proves the assertion.
\end{proof}

\appendix
\section{Verification details}
The two prime-pair counts and the tuple count use different multiplicities.
For example, the ordered vertices \((101,103,107)\) with \(s=1,R=5\)
have two anchors, \(x=97,99\). Their three half-differences are \(1,3,2\)
in edge order \(12,13,23\), with signs \(+,-,-\).
In the one-sided pair embedding any one of these pairs has exactly five
representations. These identities hold before taking a limit.

For a two-row weight an overlap already occurs for two \emph{distinct}
indices. With \(s=1,\mathcal H=\{0,2\}\) and \(v=u+1\), the four shifts
have multiplicities \(1,2,1\). Their extra sieve factor is \(\log D\).
When \(v=u\), the two distinct shifts each have multiplicity two, giving
\((\log D)^2\). Each coincidence imposes one index relation and costs
\(O(1/R)\) in counting measure. These cases check why the uniform proof
for two rough translates uses \(R/L^2\to\infty\). The graph test uses one
copy of each of the \(k-2\) auxiliary weights after both Cauchy steps,
and four weights at the opened edge, so its largest correlation has
\(j(k+2)\) shifts before merging.

The chain of analytic inputs is \eqref{eq:GT}, \eqref{eq:euler},
the collision bound, the periodic truncation, and \eqref{eq:cs} with
Davenport's bound. The limit order is \(X\to\infty\) at fixed parameters
and fixed primorial, followed by \(w\to\infty\). Fourier cancellation is
required only for the unshifted difference kernel on scale \(R\).
The principal proof uses no numerical fit or short-interval prime
asymptotic. The critical-width example also uses the ordinary prime
number theorem for the whole dyad.

\paragraph{Provenance and disclosure.}
The pair and tuple assertions in Theorem~\ref{th:weighted} were extracted
from an unpublished research record. The fixed-width prime selection in
Theorem~\ref{th:prime} was derived and checked during manuscript preparation.
The subcritical two-translate construction in Proposition~\ref{prop:roughsharp}
was recovered from the subsequent research record (Section 308DK).
The underlying mathematical research was substantially AI-generated.
OpenAI ChatGPT in Pro mode and AI review agents were used for extraction,
mathematical reconstruction, proof and prior-art reviews, and manuscript
writing. These checks are internal. Human authorship and responsibility for a final
manuscript remain pending.

\begin{thebibliography}{9}\small
\bibitem{banks} W.\ D.\ Banks and I.\ E.\ Shparlinski,
Multiple sums with the M\"obius function, \emph{Quart. J. Math.}
\textbf{77} (2026), 109--127. \href{https://doi.org/10.1093/qmath/haaf049}{doi:10.1093/qmath/haaf049}.
\href{https://arxiv.org/abs/2506.08787}{arXiv:2506.08787}.
\bibitem{dav} H.\ Davenport, On some infinite series involving arithmetical
functions (II), \emph{Quart. J. Math., Oxford Ser.} \textbf{8} (1937), 313--320.
\href{https://doi.org/10.1093/qmath/os-8.1.313}{doi:10.1093/qmath/os-8.1.313}.
\bibitem{gt} B.\ Green and T.\ Tao, Linear equations in primes,
\emph{Ann. of Math.} (2) \textbf{171} (2010), 1753--1850.
\href{https://annals.math.princeton.edu/2010/171-3/p08}{Publisher's text}.
\href{https://arxiv.org/html/math/0606088}{Readable preprint, including Appendix D}.
\bibitem{fi} J.\ Friedlander and H.\ Iwaniec, \emph{Opera de Cribro},
Amer. Math. Soc. Colloquium Publications, vol.~57 (2010), Chapter~6.
\bibitem{gtnil} B.\ Green and T.\ Tao, The M\"obius function is strongly
orthogonal to nilsequences, \emph{Ann. of Math.} (2) \textbf{175} (2012),
541--566. \href{https://annals.math.princeton.edu/2012/175-2/p03}{Publisher's text}.
\href{https://arxiv.org/html/0807.1736}{Readable preprint}.
\bibitem{lichtman} J.\ D.\ Lichtman, Averages of the M\"obius function on
shifted primes, \emph{Quart. J. Math.} \textbf{73} (2022), 729--757.
\href{https://doi.org/10.1093/qmath/haab054}{doi:10.1093/qmath/haab054}.
\href{https://arxiv.org/abs/2009.08969}{arXiv:2009.08969}.
\bibitem{rss} J.\ Rivat, A.\ S\'ark\"ozy and C.\ L.\ Stewart, Congruence
properties of the \(\Omega\)-function on sumsets, \emph{Illinois J. Math.}
\textbf{43} (1999), 1--18.
\href{https://uwaterloo.ca/pure-mathematics/sites/default/files/uploads/documents/1255985334.pdf}{Author-hosted text}.
\bibitem{shao} X.\ Shao, Narrow arithmetic progressions in the primes,
\emph{Int. Math. Res. Not. IMRN} (2017), 391--428.
\href{https://arxiv.org/abs/1509.04955}{arXiv:1509.04955}.
\bibitem{tz} T.\ Tao and T.\ Ziegler, Narrow progressions in the primes,
in \emph{Analytic Number Theory: In Honor of Helmut Maier's 60th Birthday},
Springer (2015), 357--379.
\href{https://arxiv.org/abs/1409.1327}{arXiv:1409.1327}.
\end{thebibliography}
\end{document}